\chapter*{前言}
\addcontentsline{toc}{chapter}{前言}

近年来，生成式人工智能迅速发展，新模型、新系统和新应用以惊人的速度不断涌现。然而，在这一快速变化的表象之下，许多核心思想实际上建立在一组相对精炼的数学原理之上：潜变量、似然、变分界、可逆变换、随机加噪过程、得分场、对抗式比较以及能量景观。

本书《生成式人工智能基础小册：直观数学导论》（\emph{The Little Book of Generative AI Foundations: An Intuitive Mathematical Primer}）旨在为理解这些思想提供一条紧凑而严谨的学习路径。本书并不试图罗列每一种现代架构、实现细节或最新改进，而是聚焦于使生成模型各主要家族真正可理解的数学基础。我们的目标是在不削弱数学内涵的前提下，提高这些基础知识的可达性。本书面向希望从第一性原理理解生成建模、并对数学保持好奇的研究人员、实践者和学生，可作为一部用于夯实基础的入门读物。

因此，书名中的“小册”应理解为范围上的精炼，而非深度上的浅显。本书在内容选择上有意保持克制，但在所选主题内部则力求细致严谨。核心思想按步骤展开，目的在于显露其背后的数学结构，而不是把模型当作黑箱式配方。全书关注基础而非研究前沿：已经处于前沿研究中的读者可能会熟悉许多主题，但本书希望将这些主题背后的数学逻辑解释得更加显式、直观并且连贯。

参考文献也有意保持聚焦。其主要作用是支撑正文中的关键历史脉络和技术联系，而不是把本书扩展成一份篇幅庞大的综述型书目。

\section*{写作方法与内容组织}
%
%\addcontentsline{toc}{section}{写作方法与内容组织}

本书的核心目标，是用一条连贯的数学叙事把生成模型的主要家族联系起来。章节顺序的设计，是为了让数学概念尽可能自然、紧凑地展开，而不是严格遵循历史发展顺序。各类模型按照一种便于“后续思想建立在前述思想之上”的次序引入，使读者能够看到潜变量、变分目标、扩散过程、得分场、精确密度模型、对抗学习以及基于能量的建模之间的内在联系。

全书遵循一个简单原则：只有当某种数学工具对于理解一个建模思想变得必要时，才引入它。因此，所需的线性代数、概率论、微积分、高斯代数以及密度变换工具，都嵌入在具体的建模叙事之中，而不是被单独放在前面进行冗长的预备知识回顾。

与仅做高层概览相比，本书包含了更多的数学推导。这是有意为之：生成建模中的许多思想，只有把其数学目标逐层展开之后才会真正变得清晰。推导过程按步骤给出，并在需要时引入或回顾必要概念，使数学始终与建模直觉相连接，而不是成为彼此孤立的技术模块。

本书从主成分分析（PCA）和自编码器开始，借此引入线性变换、投影、重构以及潜在结构等概念。随后，概率主成分分析（Probabilistic PCA, PPCA）把这些思想转化为本书中的第一个潜变量生成模型。

接下来进入变分生成建模。我们通过概率潜变量引入证据下界（Evidence Lower Bound, ELBO）、期望最大化（Expectation--Maximisation, EM）算法以及变分推断，并进一步扩展到变分自编码器（Variational Autoencoder, VAE）。

随后，本书转向扩散模型。首先在离散时间框架下把去噪扩散概率模型（Denoising Diffusion Probabilistic Models, DDPM）作为序列潜变量模型引入；之后介绍连续时间生成建模所需的微积分基础，并进一步转向扩散的得分视角，在这一视角下，生成过程由学习得到的得分场和逆时间随机动力学共同驱动。

接下来的一章讨论精确密度模型，包括归一化流（normalising flows）和自回归分解（autoregressive factorisations）。这些模型展示了如何通过结构设计保持似然的可计算性：要么依赖可逆变换，要么利用概率链式法则。

最后一章讨论超越可计算似然的生成建模方法。我们将介绍生成对抗网络（GAN）、Wasserstein GAN 以及基于能量的模型（Energy-Based Models, EBM），说明生成学习如何不依赖显式似然求值，而通过比较、几何差异或标量能量景观来实现。

% 两个附录汇总了高斯代数和扩散逆后验相关的辅助推导。

\section*{起源与版本说明}
%
%\addcontentsline{toc}{section}{起源与版本说明}

本书是一项由作者独立完成的学术项目，其动机来自长期以来对“如何让现代生成式人工智能的数学基础更加透明、更加易于理解”的兴趣。本书并非某项受资助科研项目、机构委托成果、正式课程开发活动，也不是针对任何特定课程或培养项目编写的官方教学文档。

本书手稿源于一个更广泛的独立研究计划，该计划关注现代生成建模的概率与数学基础。其中部分思想的早期简短版本曾以预印本形式发表，先发表于 SSRN，后经修订发表于 arXiv~\cite{chen2025plvm,chen2025unified}。本书在此基础上进一步扩展，形成一部更系统的数学导论，从经典潜变量模型一路讨论到现代生成式人工智能的基础方法。

当前手稿是本书的第二版预印本。与第一版相比，本版在全书范围内加入了更多可视化解释，并对相应正文做了少量澄清。它仍是一份工作中的预印本，后续可能继续进行修订、勘误、解释性补充或增加新的材料。未来亦可能提交经过进一步修订的版本以正式出版。

\section*{再利用说明}
%
%\addcontentsline{toc}{section}{再利用说明}

本手稿以预印本形式公开，供学术阅读、引用、研究参考和学习使用。读者可以在适当署名的前提下引用本书，并将其作为科研、学习和教学中的辅助资源。

除非适用的版权例外另有允许，或未来版本附带的其他许可另有规定，对本书内容进行大规模复制、改编、再分发，或将其转换为衍生教学、培训或课程资源，需获得作者许可。

\section*{致谢}
%
%\addcontentsline{toc}{section}{致谢}

本书源于我持续从第一性原理理解并解释现代生成式人工智能数学基础的努力。在反复尝试把线性代数、概率论、变分推断、随机建模以及现代深度生成模型中的思想连接成一条连贯叙事的过程中，许多解释逐渐变得更加清晰。

感谢更广泛的学术共同体。大量论文、书籍、讲座与讨论共同塑造了本手稿形成时所处的知识背景。文中若仍有错误、遗漏或解释不清之处，责任均在作者本人。

\vspace{1cm}

\begin{flushright}
Tianhua Chen\\
\textit{\bookdate}
\end{flushright}
